\subsection{Gradient descent and line search}
\label{subsec:gradient-descent-line-search}

With the local model from \cref{subsec:optimization-euclidean} in place, we can turn derivative information into an algorithm.
The two design questions are the search direction and the step size.
Gradient descent answers the first question by moving opposite the gradient, while fixed-step and line-search rules answer the second.

\subsubsection{Descent methods}
\label{subsec:descent-methods}

With this Euclidean setup in place, we can describe the generic update template behind many first-order algorithms.
Many iterative methods for~\eqref{eq:unconstrained-min} have the form
\begin{equation}
  x^{(k+1)} = x^{(k)} + t^{(k)} \Delta x^{(k)},
  \qquad k = 0,1,2,\dots,
  \label{eq:generic-descent}
\end{equation}
where $t^{(k)} > 0$ is a step size (learning rate) and $\Delta x^{(k)}$ is a descent direction chosen so that taking a small step decreases $f$.

\begin{definition}[Descent direction]
Let $f$ be differentiable at $x$.
A direction $d$ is a descent direction for $f$ at $x$ if $\nabla f(x)^\top d < 0$.
\end{definition}

\begin{lemma}[Why $\nabla f(x)^\top d < 0$ matters]
If $\nabla f(x)^\top d < 0$, then for all sufficiently small $t>0$,
\[
f(x+td) < f(x).
\]
\end{lemma}
\begin{proof}
Consider $\phi(t) = f(x+td)$.
Then $\phi'(0)=\nabla f(x)^\top d < 0$, so $\phi(t) < \phi(0)$ for all small enough $t>0$.
\end{proof}

\begin{algorithm}
  \caption{Generic descent method}
  \label{alg:generic-descent}
  \begin{algorithmic}[1]
    \Require Objective $f$, initial point $x^{(0)}$, rule for choosing a descent direction, rule for choosing a step size, stopping criterion
    \Ensure Final iterate $x$
    \State $k \gets 0$
    \While{stopping criterion not met}
      \State Choose $\Delta x^{(k)}$ such that $\nabla f(x^{(k)})^\top \Delta x^{(k)} < 0$.
      \State Choose a step size $t^{(k)} > 0$.
      \State $x^{(k+1)} \gets x^{(k)} + t^{(k)}\Delta x^{(k)}$.
      \State $k \gets k + 1$
    \EndWhile
    \State \Return $x^{(k)}$
  \end{algorithmic}
\end{algorithm}

\subsubsection{Step-size selection}
\label{subsec:line-search}

The remaining design choice in \cref{alg:generic-descent} is the step size.
There are two standard regimes.
If a smoothness bound is available, a fixed step size gives a clean baseline and leads to explicit rates.
If that information is unavailable or too conservative, a line search chooses a step size adaptively from local function values.

Exact line search chooses the step size by solving
\begin{equation}
  t^{(k)} = \argmin_{s>0} f\bigl(x^{(k)} + s \Delta x^{(k)}\bigr).
  \label{eq:exact-line-search}
\end{equation}
This yields the best decrease along the ray, but it is not always tractable.
More importantly, in many problems we want a step-size rule that works robustly without solving a difficult one-dimensional optimization at every iteration.

Backtracking line search (Armijo) chooses a step size by repeatedly shrinking until a \emph{sufficient decrease} condition holds.
Fix parameters $\alpha \in (0,1/2)$ and $\beta \in (0,1)$ and require the Armijo condition
\begin{equation}
  f(x+t\Delta x) \le f(x) + \alpha t \nabla f(x)^\top \Delta x,
  \label{eq:armijo-condition}
\end{equation}
to hold.
See \citet[\S9.2]{boyd_vandenberghe_2004_convex}.

\begin{algorithm}
  \caption{Backtracking line search (Armijo)}
  \label{alg:backtracking}
  \begin{algorithmic}[1]
    \Require Differentiable $f$, point $x$, descent direction $\Delta x$, parameters $\alpha \in (0,1/2)$, $\beta \in (0,1)$
    \Ensure Step size $t$
    \State $t \gets 1$
    \While{$f(x+t\Delta x) > f(x) + \alpha t \nabla f(x)^\top \Delta x$}
      \State $t \gets \beta t$
    \EndWhile
    \State \Return $t$
  \end{algorithmic}
\end{algorithm}

Typical parameter choices are:
\begin{itemize}[leftmargin=*]
  \item $\alpha$ controls how much decrease we demand; typical values are $\alpha \in (0.01,0.3)$.
  \item $\beta$ controls how aggressively we shrink; typical values are $\beta \in (0.1,0.8)$ (smaller $\beta$ means bigger reductions per backtracking step).
\end{itemize}
Backtracking is most natural in deterministic optimization where we can reliably evaluate $f(x+t\Delta x)$.
In stochastic settings (e.g.\ mini-batch training), function values are noisy, so line search is less common; one typically uses fixed or scheduled learning rates, possibly with averaging or momentum.
For the running logistic regression example in \cref{ex:logistic-reg}, this means that once we choose a direction $\Delta\theta$, we can evaluate the full objective $s(\theta+t\Delta\theta)$ exactly and shrink $t$ until the sufficient-decrease test passes.

\begin{figure}[t]
  \centering
  \begin{tikzpicture}
  % --- Main plot (geometry of Armijo condition) ---
  \begin{axis}[
    name=main,
    width=0.95\linewidth,
    height=6.2cm,
    xmin=0, xmax=1.25,
    ymin=0.0, ymax=0.8,
    axis lines=left,
    ticks=none,
    xlabel={$t$},
    ylabel={$\,\phi(t)=f(x+t\Delta x)$},
    clip=false,
  ]
    % A representative 1D slice along the search direction.
    \addplot[black,thick,domain=0:1.25,samples=220] {(x-0.4)^2+0.3};
    % Tangent at t=0:  phi(0) + t phi'(0), shown only near t=0.
    \addplot[stat221Blue,thick,domain=0:0.55,samples=2] {0.46-0.8*x};
    % Armijo (sufficient decrease) line: phi(0) + alpha t phi'(0), alpha=0.2.
    \addplot[stat221Purple,thick,domain=0:1.25,samples=2] {0.46-0.16*x};

    % Mark t0 where phi(t) intersects the Armijo line.
    \addplot[only marks,mark=*,mark size=1.6pt] coordinates {(0.64,0.3576)};
    \addplot[densely dashed,stat221Gray] coordinates {(0.64,0) (0.64,0.3576)};
    \node[anchor=north] at (axis cs:0.64,0) {$t_0$};

    % Show the typical backtracking output t = beta^k (some value <= t0).
    \addplot[densely dotted,stat221Purple] coordinates {(0.49,0) (0.49,0.70)};
    \node[anchor=north,stat221Purple] at (axis cs:0.49,0) {$\beta^k$};

    % Also mark the initial trial step t=1.
    \addplot[densely dashed,stat221Gray] coordinates {(1,0) (1,0.78)};
    \node[anchor=north] at (axis cs:1,0) {$1$};

    % Labels (kept lightweight; the caption does the heavy lifting).
    \node[black,anchor=west] at (axis cs:1.02,0.68) {$\phi(t)$};
    \node[stat221Purple,anchor=west] at (axis cs:0.82,0.44) {$\phi(0)+\alpha t\,\phi'(0)$};
    \node[stat221Blue,anchor=west] at (axis cs:0.12,0.11) {$\phi(0)+t\,\phi'(0)$};
  \end{axis}

  % --- Case plots (what the line search returns) ---
  \begin{axis}[
    at={(main.south west)}, anchor=north west,
    yshift=-0.9cm,
    width=0.46\linewidth,
    height=3.3cm,
    xmin=0, xmax=1.1,
    ymin=0.0, ymax=0.55,
    axis lines=left,
    ticks=none,
    xlabel={$t$},
    ylabel={},
    title={Case 1: $t_0>1$},
    title style={font=\small, yshift=-0.35em},
  ]
    \addplot[black,thick,domain=0:1.1,samples=120] {0.46-0.5*x+0.3*x^2};
    \addplot[stat221Purple,thick,domain=0:1.1,samples=2] {0.46-0.1*x};
    \addplot[densely dashed,stat221Gray] coordinates {(1,0) (1,0.52)};
    \addplot[only marks,mark=*,mark size=1.3pt] coordinates {(1,0.26)};
    \node[anchor=west,font=\small] at (axis cs:0.02,0.50) {output: $t=1$};
  \end{axis}

  \begin{axis}[
    at={(main.south west)}, anchor=north west,
    xshift=0.49\linewidth,
    yshift=-0.9cm,
    width=0.46\linewidth,
    height=3.3cm,
    xmin=0, xmax=1.1,
    ymin=0.0, ymax=0.80,
    axis lines=left,
    ticks=none,
    xlabel={$t$},
    ylabel={},
    title={Case 2: $t_0<1$},
    title style={font=\small, yshift=-0.35em},
  ]
    \addplot[black,thick,domain=0:1.1,samples=160] {(x-0.4)^2+0.3};
    \addplot[stat221Purple,thick,domain=0:1.1,samples=2] {0.46-0.16*x};
    \addplot[densely dotted,stat221Purple] coordinates {(0.49,0) (0.49,0.78)};
    \addplot[only marks,mark=*,mark size=1.3pt,stat221Purple] coordinates {(0.49,0.3081)};
    \node[anchor=west,font=\small] at (axis cs:0.02,0.72) {output: $t=\beta^k$};
  \end{axis}
\end{tikzpicture}


  \caption{Geometry of the Armijo (sufficient decrease) condition for $\phi(t)=f(x+t\Delta x)$.
  The Armijo line $\phi(0)+\alpha t\,\phi'(0)$ intersects $\phi$ at $t_0$.
  Backtracking starts at $t=1$ and returns $t=\beta^k$ (possibly $t=1$) such that $\phi(t)\le \phi(0)+\alpha t\,\phi'(0)$.
  }
  \label{fig:armijo-backtracking}
\end{figure}

Backtracking line search is a cheap way to find a step size that is ``small enough'' without knowing global smoothness constants.
The Armijo condition compares the actual decrease $f(x+t\Delta x)-f(x)$ to the first-order prediction $t\nabla f(x)^\top \Delta x$.
If $f$ is reasonably well-approximated by its tangent for small steps, repeated shrinking will eventually make the inequality hold (see \citet[\S9.2]{boyd_vandenberghe_2004_convex}).


When that scale is unknown, backtracking gives a local replacement.
We focus on gradient descent with a backtracking line search (see \citet[\S9.3]{boyd_vandenberghe_2004_convex}).
\Cref{alg:gd-backtracking} is the specialization of \cref{alg:generic-descent} with direction $\Delta x^{(k)}=-\nabla f(x^{(k)})$ and step size chosen by \cref{alg:backtracking}.
The analysis is modular: smoothness implies that \emph{some} step sizes satisfy Armijo (\cref{lem:armijo-sufficient}), which implies backtracking never shrinks ``too much'' (\cref{prop:bt-lower-bound}); strong convexity then links the gradient norm to suboptimality (\cref{cor:sc-links}), yielding a linear rate (\cref{thm:gd-linear}).
In this subsection we work with unconstrained objectives $f:\R^d\to\R$, so every trial point $x-t\nabla f(x)$ is feasible.

\begin{algorithm}
  \caption{Gradient descent with backtracking line search}
  \label{alg:gd-backtracking}
  \begin{algorithmic}[1]
    \Require Differentiable $f:\R^d\to\R$, initial point $x^{(0)}$, parameters $\alpha \in (0,1/2)$, $\beta \in (0,1)$, stopping criterion
    \Ensure Final iterate $x$
    \State $k \gets 0$
    \While{stopping criterion not met}
      \State $g^{(k)} \gets \nabla f(x^{(k)})$
      \State $\Delta x^{(k)} \gets -g^{(k)}$
      \State $t^{(k)} \gets$ \Call{BacktrackingLineSearch}{$f, x^{(k)}, \Delta x^{(k)}, \alpha, \beta$} \Comment{\cref{alg:backtracking}}
      \State $x^{(k+1)} \gets x^{(k)} + t^{(k)}\Delta x^{(k)}$
      \State $k \gets k + 1$
    \EndWhile
    \State \Return $x^{(k)}$
  \end{algorithmic}
\end{algorithm}

The update has the form
\begin{equation}
  x^{(k+1)} = x^{(k)} - t^{(k)} \nabla f(x^{(k)}),
  \label{eq:gd-update}
\end{equation}
where $t^{(k)}$ is the step size returned by backtracking.

\begin{lemma}[A simple sufficient step size for the Armijo condition]
\label{lem:armijo-sufficient}
Assume $f:\R^d\to\R$ is $M$-smooth and take $\alpha \in (0,1/2)$.
Then for any $x$ and any step size $t \in (0,1/M]$,
\begin{equation}
  f(x-t\nabla f(x)) \le f(x) - \alpha t \norm{\nabla f(x)}_2^2.
  \label{eq:armijo-sufficient}
\end{equation}
\end{lemma}
\begin{proof}
By the upper bound in \cref{thm:quadratic-bounds} (the ``descent lemma''), with $y=x-t\nabla f(x)$,
\[
f(x-tg) \le f(x) + \nabla f(x)^\top(-tg) + \frac{M}{2}\norm{-tg}_2^2
= f(x) - t\norm{g}_2^2 + \frac{M t^2}{2}\norm{g}_2^2,
\]
where $g=\nabla f(x)$.
If $t \le 1/M$ then $-t + (Mt^2/2) \le -t/2$, so
\[
f(x-tg) \le f(x) - \frac{t}{2}\norm{g}_2^2 \le f(x) - \alpha t \norm{g}_2^2,
\]
since $\alpha<1/2$.
\end{proof}

\begin{proposition}[Backtracking returns a step size bounded away from zero]
\label{prop:bt-lower-bound}
Assume $f:\R^d\to\R$ is $M$-smooth and backtracking is run with $\alpha \in (0,1/2)$ and $\beta \in (0,1)$.
Then the returned step size $t$ satisfies
\begin{equation}
  t \ge \min\left\{1,\frac{\beta}{M}\right\}.
  \label{eq:bt-lower-bound}
\end{equation}
\end{proposition}
\begin{proof}
If backtracking accepts $t=1$, then \eqref{eq:bt-lower-bound} holds.
Otherwise, at least one shrink occurs, so the final accepted $t$ has the form $t=\beta^j$ with $j\ge 1$ and the previous trial value $t/\beta$ was rejected.
By \cref{lem:armijo-sufficient}, any $s\le 1/M$ satisfies the Armijo condition, so a rejected step must satisfy $t/\beta > 1/M$.
Thus $t > \beta/M$.
Since this case only occurs when $t<1$, we also have $\beta/M < 1$, so $t \ge \min\{1,\beta/M\}$.
\end{proof}

\begin{theorem}[Linear convergence of GD with backtracking for strongly convex and smooth $f$]
\label{thm:gd-linear}
Assume $f:\R^d\to\R$ is twice differentiable and satisfies
\[
  m I \preceq \Hess f(x) \preceq M I \qquad \forall x \in \R^d,
\]
with $0<m\le M$.
Let $(x^{(k)})$ be the gradient descent iterates \eqref{eq:gd-update} with backtracking parameters $\alpha \in (0,1/2)$ and $\beta \in (0,1)$, and let $x^\star=\argmin_{x\in\R^d} f(x)$.
Then
\begin{equation}
  f(x^{(k)}) - f(x^\star)
  \le c^k \bigl(f(x^{(0)}) - f(x^\star)\bigr),
  \qquad
  c = 1 - \min\left\{2\alpha m, \frac{2\alpha\beta m}{M}\right\}.
  \label{eq:gd-linear}
\end{equation}
\end{theorem}
\begin{proof}
By the Armijo condition with $\Delta x^{(k)}=-\nabla f(x^{(k)})$,
\[
f(x^{(k+1)}) \le f(x^{(k)}) - \alpha t^{(k)} \norm{\nabla f(x^{(k)})}_2^2.
\]
By \cref{prop:bt-lower-bound}, $t^{(k)} \ge \min\{1,\beta/M\}$.
By \cref{cor:sc-links}(b), $\norm{\nabla f(x^{(k)})}_2^2 \ge 2m\bigl(f(x^{(k)})-f(x^\star)\bigr)$.
Combining,
\[
f(x^{(k+1)})-f(x^\star)
\le \left(1 - 2\alpha m \min\left\{1,\frac{\beta}{M}\right\}\right)\bigl(f(x^{(k)})-f(x^\star)\bigr).
\]
Since $2\alpha m \min\{1,\beta/M\} = \min\{2\alpha m, 2\alpha\beta m/M\}$, this yields \eqref{eq:gd-linear} by iteration.
\end{proof}

It is useful to express the rate in terms of the condition number $\kappa=M/m$.
When $\kappa$ is large (poor conditioning), the contraction factor in \eqref{eq:gd-linear} is approximately
\[
c \approx 1 - \frac{2\alpha\beta}{\kappa},
\]
so convergence can be slow.
Geometrically, the level sets are long narrow valleys and vanilla gradient descent tends to ``zig-zag'' (see \citet[\S9.3 and Ex.~9.3.2]{boyd_vandenberghe_2004_convex}).
This motivates changing the geometry of the update via preconditioning.


\subsubsection{Gradient descent on quadratics}

With step-size rules in hand, quadratics provide the cleanest exact benchmark for how first-order methods behave under conditioning.
When a smoothness scale is available in advance, constant-step gradient descent gives a simple baseline, and for strongly convex quadratics the rate can be computed exactly from the spectrum.
If we know a smoothness constant $M$ in advance, we can choose a constant step size without a line search.
Consider the update
\begin{equation}
  x^{(k+1)} = x^{(k)} - t\,\nabla f(x^{(k)}),
  \qquad t>0.
  \label{eq:gd-constant-stepsize}
\end{equation}

\begin{theorem}[Sublinear convergence for smooth convex objectives]
\label{thm:gd-sublinear}
Assume $f$ is convex and $M$-smooth, and let $x^\star \in \argmin f$.
If we run \eqref{eq:gd-constant-stepsize} with step size $t=1/M$, then for every $k\ge 1$,
\begin{equation}
  f(x^{(k)}) - f(x^\star) \le \frac{M}{2k}\norm{x^{(0)}-x^\star}_2^2.
  \label{eq:gd-sublinear}
\end{equation}
\end{theorem}
\begin{proof}
Write $x_k=x^{(k)}$ and $g_k=\nabla f(x_k)$.
By the smoothness upper bound in \cref{thm:quadratic-bounds} with $y=x_k-tg_k$,
\[
  f(x_{k+1})
  \le f(x_k) - t\norm{g_k}_2^2 + \frac{Mt^2}{2}\norm{g_k}_2^2
  = f(x_k) - \Bigl(t-\frac{Mt^2}{2}\Bigr)\norm{g_k}_2^2.
\]
With $t=1/M$ this becomes
\begin{equation}
  f(x_k)-f(x_{k+1}) \ge \frac{1}{2M}\norm{g_k}_2^2.
  \label{eq:gd-one-step-decrease}
\end{equation}

Next expand the squared distance to $x^\star$:
\[
  \norm{x_{k+1}-x^\star}_2^2
  = \norm{x_k-x^\star}_2^2 - 2t\,g_k^\top(x_k-x^\star) + t^2\norm{g_k}_2^2.
\]
By convexity, $g_k^\top(x_k-x^\star)\ge f(x_k)-f(x^\star)$.
Also \eqref{eq:gd-one-step-decrease} implies $t^2\norm{g_k}_2^2 \le 2t\,(f(x_k)-f(x_{k+1}))$.
Combine these two bounds to get
\[
  \norm{x_{k+1}-x^\star}_2^2
  \le \norm{x_k-x^\star}_2^2 - 2t\,(f(x_k)-f(x^\star)) + 2t\,(f(x_k)-f(x_{k+1}))
  = \norm{x_k-x^\star}_2^2 - 2t\,(f(x_{k+1})-f(x^\star)).
\]
Rearranging,
\[
  f(x_{k+1})-f(x^\star)
  \le \frac{1}{2t}\Bigl(\norm{x_k-x^\star}_2^2 - \norm{x_{k+1}-x^\star}_2^2\Bigr).
\]
Sum from $k=0$ to $K-1$ to obtain
\[
  \sum_{k=0}^{K-1}\bigl(f(x_{k+1})-f(x^\star)\bigr)
  \le \frac{1}{2t}\norm{x_0-x^\star}_2^2.
\]
Finally, \eqref{eq:gd-one-step-decrease} implies $f(x_{k+1})\le f(x_k)$, so
\[
  K\bigl(f(x_K)-f(x^\star)\bigr)
  \le \sum_{k=0}^{K-1}\bigl(f(x_{k+1})-f(x^\star)\bigr)
  \le \frac{1}{2t}\norm{x_0-x^\star}_2^2.
\]
With $t=1/M$ this yields \eqref{eq:gd-sublinear}.
\end{proof}

The rate \eqref{eq:gd-sublinear} is \emph{sublinear} in $k$: to drive the suboptimality below $\varepsilon$ we need $k=O(1/\varepsilon)$ iterations.

\begin{corollary}[Linear convergence under strong convexity with a fixed step size]
\label{cor:gd-fixed-linear}
If, in addition, $f$ is $m$-strongly convex, then with step size $t=1/M$,
\begin{equation}
  f(x^{(k)})-f(x^\star)
  \le \left(1-\frac{m}{M}\right)^k\bigl(f(x^{(0)})-f(x^\star)\bigr).
  \label{eq:gd-fixed-linear}
\end{equation}
\end{corollary}
\begin{proof}
With $t=1/M$, \eqref{eq:gd-one-step-decrease} gives
\[
  f(x_{k+1}) \le f(x_k) - \frac{1}{2M}\norm{\nabla f(x_k)}_2^2.
\]
By \cref{cor:sc-links}(b), $\norm{\nabla f(x_k)}_2^2 \ge 2m(f(x_k)-f(x^\star))$.
Combine and rearrange to get
\[
  f(x_{k+1})-f(x^\star) \le \left(1-\frac{m}{M}\right)\bigl(f(x_k)-f(x^\star)\bigr),
\]
then iterate.
\end{proof}

Geometric rates of the form $c^k$ are often referred to as \emph{linear convergence} (since $\log(f(x^{(k)})-f(x^\star))$ decays roughly linearly in $k$).

For general smooth strongly convex objectives, the bound in \cref{cor:gd-fixed-linear} already shows a geometric rate.
For quadratics we can compute the rate exactly via eigenvalues, and in particular identify the optimal constant step size and its dependence on the condition number.

\begin{proposition}[Gradient descent on a strongly convex quadratic]
\label{prop:gd-quadratic}
Let $A\in\R^{d\times d}$ be symmetric positive definite with eigenvalues in $[m,M]$, and let
\[
f(x)=\frac12 x^\top A x - b^\top x.
\]
Run gradient descent with a constant step size $t>0$,
\[
x_{k+1}=x_k - t\nabla f(x_k)=x_k - t(Ax_k-b).
\]
Then $x^\star=A^{-1}b$ and the error $e_k\coloneqq x_k-x^\star$ satisfies
\[
e_{k+1}=(I-tA)e_k,
\qquad
e_k=(I-tA)^k e_0.
\]
Moreover, if $0<t<2/M$ then
\[
\norm{e_k}_2 \le \rho(t)^k \norm{e_0}_2,
\qquad
\rho(t) \coloneqq \max\{|1-tm|,|1-tM|\}\in(0,1),
\]
and the function gap contracts as
\[
f(x_k)-f(x^\star)=\frac12\norm{e_k}_A^2 \le \frac{M}{2}\rho(t)^{2k}\norm{e_0}_2^2.
\]
The choice $t^\star=\frac{2}{m+M}$ minimizes $\rho(t)$, yielding
\[
\rho(t^\star)=\frac{M-m}{M+m} = \frac{\kappa-1}{\kappa+1},
\qquad \kappa\coloneqq M/m.
\]
\end{proposition}
\begin{proof}
Since $\nabla f(x)=Ax-b$ and $x^\star=A^{-1}b$, we have
\[
e_{k+1}=x_{k+1}-x^\star = x_k-x^\star - t(Ax_k-b) = e_k - tAe_k = (I-tA)e_k,
\]
and iterating gives $e_k=(I-tA)^k e_0$.

Diagonalize $A=Q\Lambda Q^\top$ with $Q$ orthogonal and $\Lambda=\mathrm{diag}(\lambda_1,\dots,\lambda_d)$.
Then $I-tA = Q(I-t\Lambda)Q^\top$, so
\[
\norm{e_k}_2 = \norm{(I-tA)^k e_0}_2
\le \norm{(I-tA)^k}_2\,\norm{e_0}_2
= \max_i |1-t\lambda_i|^k \norm{e_0}_2.
\]
The map $\lambda\mapsto |1-t\lambda|$ is convex, so its maximum over $\lambda\in[m,M]$ occurs at an endpoint, giving
$\max_i |1-t\lambda_i| = \max\{|1-tm|,|1-tM|\} = \rho(t)$.
If $0<t<2/M$ then $|1-t\lambda_i|<1$ for every $\lambda_i\in[m,M]$, hence $\rho(t)\in(0,1)$.

Finally, completing the square gives
$f(x)-f(x^\star)=\frac12 (x-x^\star)^\top A (x-x^\star)=\frac12\norm{e}_A^2$,
and since $A\preceq MI$ we have $\norm{e}_A^2 \le M\norm{e}_2^2$.
The minimizer of $\rho(t)=\max\{|1-tm|,|1-tM|\}$ occurs when the two absolute values are equal, i.e.\ when $1-tm=-(1-tM)$.
This gives $t^\star=2/(m+M)$ and $\rho(t^\star)=(M-m)/(M+m)$.
\end{proof}

When $M$ is unknown (or too conservative), a line search such as backtracking is a robust way to pick step sizes automatically; we analyze this next.

